\chapter{绪论}

\section{第一节：什么是信号}
信号是携带信息的物理量，是信息的载体，如电信号、光信号、声信号等，广泛存在于通信、控制、检测等工程领域。

\section{第二节：信号的描述与分类}
\subsection{一、信号的数学描述}
\begin{enumerate}
\item 数学表达式描述
  \begin{itemize}
  \item 核心：用函数关系 $f(t)$ / $f(n)$ 表征信号随时间/离散变量的变化规律
  \item 示例：正弦信号 $f(t) = A\sin(\omega t + \varphi)$、单位阶跃信号 $u(t)$
  \end{itemize}
\item 形态描述
  \begin{itemize}
  \item 核心：通过波形图、时域曲线等可视化方式描述信号特征
  \item 关键维度：幅值、周期、上升/下降沿、突变点等
  \end{itemize}
\end{enumerate}

\subsection{二、信号的分类}
\begin{enumerate}
\item 按信号中独立变量的个数
  \begin{itemize}
  \item 一维信号：仅含一个独立变量（如时间 $t$），如语音信号 $f(t)$
  \item 多维信号：含多个独立变量，如图像信号 $f(x,y)$、视频信号 $f(x,y,t)$
  \end{itemize}
\item 按信号的时间特性
  \begin{itemize}
  \item 确定信号：取值可由确定函数描述（如正弦信号、方波信号）
  \item 随机信号：取值无确定规律，仅能统计描述（如噪声、地震波）
  \end{itemize}
\item 按时间函数特性
  \begin{itemize}
  \item 连续时间信号：在连续时间范围内有定义（$t \in \mathbb{R}$），记为 $f(t)$
    \begin{itemize}
    \item 数学描述：$f(t), t \in (-\infty, +\infty)$，幅值可连续/离散（模拟信号是幅值也连续的连续时间信号）
    \end{itemize}
  \item 离散时间信号：仅在离散时间点有定义（$t = nT, n \in \mathbb{Z}$），记为 $f(n)$
    \begin{itemize}
    \item 数学描述：$f(n), n = 0, \pm1, \pm2, \dots$，也称为序列
    \end{itemize}
  \end{itemize}
\item 按周期性
  \begin{itemize}
  \item 周期信号：存在非零常数 $T$，满足 $f(t+T) = f(t)$（连续）/ $f(n+N) = f(n)$（离散）
    \begin{itemize}
    \item 判定规则：多个周期信号叠加，频率比为有理数时仍为周期信号
    \end{itemize}
  \item 非周期信号：无固定周期，如脉冲信号、衰减指数信号
  \end{itemize}
\item 按能量特点
  \subsubsection{信号的能量与功率分类}
  \begin{quote}
  以下定义基于 \textbf{1Ω 电阻} 负载，瞬时功率 $P = |f(t)|^2$
  \end{quote}
  
  \paragraph{1. 信号能量定义}
  \[
  E = \frac{1}{R} \int_{-\infty}^{\infty} |f(t)|^2 dt = \lim_{T \to \infty} \int_{-T}^{T} |f(t)|^2 dt
  \]
  \begin{quote}
  【说明】：$E$ 为信号总能量，单位为焦耳（J）
  \end{quote}
  
  \paragraph{2. 信号平均功率定义}
  \[
  P = \lim_{T \to \infty} \frac{1}{2T} \int_{-T}^{T} |f(t)|^2 dt = \lim_{T \to \infty} \frac{1}{T} \int_{-\frac{T}{2}}^{\frac{T}{2}} |f(t)|^2 dt
  \]
  \begin{quote}
  【说明】：$P$ 为信号平均功率，单位为瓦特（W）
  \end{quote}
  
  \paragraph{信号分类规则}
  根据信号的能量 $E$ 和功率 $P$ 取值，信号分为三类：
  \begin{itemize}
  \item \textbf{能量信号}
    \begin{quote}
    判定条件：$E < \infty$（能量有限），且 $P = 0$\\
    示例：矩形脉冲、衰减指数信号等\\
    【特点】：信号能量集中，随时间快速衰减，总能量有限
    \end{quote}
  \item \textbf{功率信号}
    \begin{quote}
    判定条件：$P < \infty$（功率有限），且 $E = \infty$\\
    示例：正弦信号、方波信号等周期信号\\
    【特点】：持续存在，功率稳定，但总能量无限
    \end{quote}
  \item \textbf{非能量非功率信号}
    \begin{quote}
    判定条件：$E = \infty$ 且 $P = \infty$\\
    示例：随时间无界增长的信号（如 $f(t) = t$）\\
    【特点】：既无有限总能量，也无稳定平均功率
    \end{quote}
  \end{itemize}
\end{enumerate}

\section{第三节：信号的分析与处理}
核心分析方法包括时域分析、频域分析、时频分析；处理手段包括滤波、降噪、特征提取、调制/解调等。

\section{第四节：自动控制系统中的信号分析与处理}
核心应用场景：
\begin{enumerate}
\item 控制系统中的信号类型：输入信号、反馈信号、误差信号
\item 分析方法：傅里叶变换分析系统频率响应、Z变换分析离散控制系统
\item 处理应用：噪声滤除、信号采样与重构、故障信号识别等
\end{enumerate}